\section{为什么是深度学习？为什么是现在？}

\subsection{传统机器学习的局限}

传统机器学习依赖于\textbf{手工设计的特征}（hand-engineered features）。以人脸检测为例，工程师需要手动定义如何从图像中提取边缘、曲线、五官等特征。这种方法存在三个根本问题：

\begin{enumerate}
\item \textbf{耗时}：特征工程需要大量领域专家的时间和经验
\item \textbf{脆弱}：手工特征对环境变化（光照、角度等）不具鲁棒性
\item \textbf{不可扩展}：每个新任务都需要重新设计特征集
\end{enumerate}

\begin{importantbox}{深度学习的核心思想}
深度学习的关键突破在于：\textbf{让模型自动从数据中学习特征的层次化表示}（hierarchical feature representations）。低层特征（如边缘和线条）自动组合成中层特征（如眼睛、鼻子），再进一步组合成高层特征（如面部结构）。无需人工干预。
\end{importantbox}

\begin{figure}[H]
\centering
\includegraphics[width=0.85\textwidth]{slides-images/slide-018.jpg}
\caption{特征的层次化学习：从低级特征到高级特征\protect\footnotemark}
\end{figure}
\footnotetext{Slides 第 18 页。}

\subsection{三大驱动力}

尽管神经网络的理论可以追溯到数十年前（1952 年的随机梯度下降、1958 年的感知机、1986 年的反向传播、1995 年的深度卷积网络），但其真正的爆发得益于三大驱动力的同时成熟：

\begin{figure}[H]
\centering
\includegraphics[width=0.85\textwidth]{slides-images/slide-019.jpg}
\caption{神经网络的历史与深度学习爆发的三大驱动力\protect\footnotemark}
\end{figure}
\footnotetext{Slides 第 19 页。}

\begin{table}[H]
\centering
\begin{tabular}{lp{10cm}}
\toprule
\textbf{驱动力} & \textbf{具体表现} \\
\midrule
Big Data & 大规模数据集（如 ImageNet、Wikipedia）的出现，数据收集与存储成本大幅下降 \\
Hardware & GPU 的大规模并行计算能力使得训练深度网络成为可能 \\
Software & 开源框架（TensorFlow、PyTorch、Keras）使得构建和训练深度学习模型变得简便 \\
\bottomrule
\end{tabular}
\caption{深度学习爆发的三大驱动力}
\end{table}

\begin{knowledgebox}{历史视角：这些技术并不新}
Amini 特别强调，本讲中介绍的几乎所有技术——感知机、反向传播、梯度下降——都是\textbf{数十年前}发明的。真正新的是它们在大数据、强算力和好工具三者交汇下所展现出的惊人能力。理解这些基础对于推动领域进一步发展至关重要。
\end{knowledgebox}

\subsection{本章小结}

深度学习相比传统机器学习的核心优势在于自动学习特征表示，无需人工特征工程。尽管基础理论已有数十年历史，但大数据、GPU 硬件和开源软件的三重驱动力使得深度学习在近年来迎来爆发式发展。

